\Askhsh{22244}{1}
\begin{ylh}
    \item Γεωμετρία
\end{ylh}

Ένας τετράγωνος κήπος ΑΒΓΔ έχει πλευρά 10m.

\begin{alist}
    \item Για την άρδευση του κήπου τοποθετούνται μηχανισμοί ποτίσματος στις απέναντι κορυφές Α και Γ. Ο κάθε μηχανισμός ποτίζει έναν κυκλικό τομέα γωνίας $90^\circ$ και ακτίνας 10m.
    \begin{rlist}
        \item Να βρείτε το εμβαδόν της επιφάνειας του κήπου που ποτίζει ο κάθε μηχανισμός. \monades{5}
        \item Αν οι δύο μηχανισμοί ποτίζουν ολόκληρο τον κήπο, να βρείτε το εμβαδόν της επιφάνειας του κήπου που ποτίζεται και από τους δύο μηχανισμούς. \monades{7}
    \end{rlist}
    
    \item Για την άρδευση του κήπου τοποθετούνται τέσσερις μηχανισμοί ποτίσματος στις κορυφές Α, Β, Γ και Δ. Ο κάθε μηχανισμός ποτίζει έναν κυκλικό τομέα γωνίας $90^\circ$ και ακτίνας 5m, όπως φαίνεται στο παρακάτω σχήμα.
    
    \begin{wrapfigure}{r}{5.5cm}
    \centering
    \begin{tikzpicture}[scale=0.9]
        \usetikzlibrary{patterns}
        % Ορισμός κορυφών τετραγώνου
        \tkzDefPoint(0,0){D}
        \tkzDefPoint(5,0){C}
        \tkzDefPoint(5,5){B}
        \tkzDefPoint(0,5){A}
        
        % Χρήση scope και clip για να περιοριστεί η σχεδίαση εντός του τετραγώνου
        \begin{scope}
            \tkzClipPolygon(A,B,C,D)
            % Ορίζουμε μια διαδρομή που είναι η ένωση των τεσσάρων κύκλων και τη γεμίζουμε.
            % Η ακτίνα είναι η μισή πλευρά (5/2 = 2.5).
            \path[pattern=crosshatch, pattern color=maincolor!60] 
                (A) circle (2.5)
                (B) circle (2.5)
                (C) circle (2.5)
                (D) circle (2.5);
        \end{scope}
        
        % Σχεδίαση του περιγράμματος του τετραγώνου
        \tkzDrawPolygon[thick](A,B,C,D)
        
        % Ετικέτες στις κορυφές
        \tkzLabelPoint[above left](A){A}
        \tkzLabelPoint[above right](B){B}
        \tkzLabelPoint[below right](C){$\Gamma$}
        \tkzLabelPoint[below left](D){$\Delta$}
    \end{tikzpicture}
    \end{wrapfigure}
    
    \begin{rlist}
        \item Να βρείτε το εμβαδόν της επιφάνειας του κήπου που ποτίζει ο κάθε μηχανισμός. \monades{5}
        \item Να βρείτε το εμβαδόν της επιφάνειας του κήπου που δεν ποτίζεται. \monades{8}
    \end{rlist}
    
\end{alist}